\chapter{TTC – Telecomunicaciones}

\section{Módulo de Rastreo, Telemetría y Comandos (TTC)}
El subsistema TTC (Telemetry, Tracking and Command) establece el enlace unificado de RF entre el Segmento Tierra (Estación INRAS y redes afiliadas) y el Segmento Espacial intetizado por el OBC.

\section{Arquitectura de Transmisión LEO}

\subsection{Enlaces UHF/VHF y Segmentación}
El sistema de modulación interno asume todas las transmisiones segmentadas bajo sub-sistemas radiales acoplados:
\begin{itemize}
    \item \textbf{Beacon (Baliza de Vida):} Transmisor independiente periódicamente encendido por hardware que reporta telemetría vital comprimida en CW Morse, AFSK o GMSK. Es la última barrera indicando modo de vida o reencendido sin requerir computación completa.
    \item \textbf{Enlace de Subida / Bajada (Uplink / Downlink UHF/VHF):} Tránsito bidireccional semidúplex o dúplex. El Payload transfiere hacia Downlink archivos crudos segmentados mediante paquetes CCSDS por el bus SPI al módem. El Uplink recibe comandos correctivos de tierra verificando CRC y un código seguro (HMAC/SHA) garantizando autorización para alterar la memoria operativa de apuntamiento.
\end{itemize}

\section{Antenas y Mecanismos de Despliegue Eléctrico}
\subsection{Arquitectura Fija y Retención Balística}
El diseño de antena (comúnmente dipolos/monopolos o cintas metálicas cruzadas para ganancia isotrópica) requiere plegado volumétrico absoluto ajustado al \textit{Form Factor} de 2U.
\begin{itemize}
    \item Se anclan inicialmente perimetralmente con alambres pre-tensionados mediante un polímero de nylon y unidos físicamente mediante resistencias quemadoras adyacentes a la estructura de PCB externa.
\end{itemize}

\subsection{Secuencia de Encendido Activo}
La responsabilidad total de la liberación subyace al \textit{Modo de Inicio} proveniente de la computadora de a bordo (OBC). Superada la ventana neutral de 30 a 45 min, el sistema eléctrico inyecta altas corrientes por una franja temporal (aprox. 5 a 10 segundos) originando la desgasificación y termofusión del hilo Nylon y permitiendo el despliegue dinámico por rebote natural a la posición ortogonal predispuesta de cobertura de radio LEO omnidireccional.
